\scanpage{17}
\mspage{8}
% archive sheet number: 17

la relation de départ en se restreignant à $\unsure{\mathscr L|L}$.

Le cas ``quasi-ensembliste'' est celui où on suppose

\uline{\unsure{Mog|Mag} quasi}\quad Pour $x,y\in\mathscr L$, $x\neq y\Longrightarrow x\parallel y$.

\uline{Proposition}-\uline{\unsure{Scholie|Scolie}} \unsure{\deleted{Les}|Les} \inserted{structures de} magmas \deleted{quasi}\inserted{pré}-ensemblistes sur $M$ correspondent aux couples $(\leq,\ll)$ de deux relations d'ordre sur $M$ sans plus, satisfaisant les seules conditions

\uline{Mog 1}\quad $X\leq Y\Longrightarrow X\ll Y$

\uline{Mog 2}\quad $\forall\,X\ll Y$, l'ens. $\widetilde Y^X=\{Y'\in\widetilde Y\mid X\ll Y'\}$ a un plus petit élément

\uline{Mog $L\deleted{\unsure{3}}$ 1}\quad $\forall\,X\in M,\ \operatorname{omb}(X)^\circ\neq\varnothing$, i.e. $\exists\,x\in\mathscr L$, \inserted{avec} $x\overset{\circ}{\ll}X$ i.e. $\deleted{x}\widetilde X^x=\{X\}$.

Voir exemple dans CF $\underline{\overline{\mathrm{VIII}}}$ p.57 pour montrer que ceci n'implique pas que $M$ se \uline{plonge} dans $\operatorname{Figel}(\mathscr L)$.

\deleted{Les magmas \unsure{non|pré}-ensemblistes \unsure{si on a|où on a}}

Pour mémoire : si $M$ satisfait la condition supplémentaire Mog $L$ 1 (sans être même quasi ens --- ce n'a aucune importance \deleted{\unsure{?}} ici) alors on a pour $F\in\unsure{\mathfrak f|\mathscr F|\mathcal F|f}$

\[
\begin{aligned}
\text{a)}\qquad\unsure{\mathrm{Ssfig}|\mathrm{Ssfi}_0}(F)&\xrightarrow{\sim}\mathrm{Ssfig}(\operatorname{mulomb}(F))\\
F'&\longmapsto\operatorname{mulomb}(F')
\end{aligned}
\]
(\unsure{iso.|isom.} d'ens ordonnés). Par contre, même dans le cas quasi-ensembliste, on \deleted{peut}\inserted{n'a} pas nécessairement :
